\documentclass[11pt,a4paper]{article}
\usepackage[utf8]{inputenc}
\usepackage[T1]{fontenc}
\usepackage[margin=2.5cm]{geometry}
\usepackage{amsmath,amssymb,amsthm}
\usepackage{mathtools}
\usepackage{graphicx}
\usepackage{booktabs}
\usepackage{hyperref}
\usepackage{listings}
\usepackage{enumitem}
\usepackage{xcolor}

\hypersetup{
    colorlinks=true,
    linkcolor=blue!60!black,
    citecolor=blue!60!black,
    urlcolor=blue!60!black
}

\title{\textsf{Oxydome}\\[0.3em]
       \large An Ethereum-like Account and UTXO Hybrid Blockchain with\\
       \large Tenant Chains, Fuel Currencies, and Authenticated Auditing}
\author{
  \texttt{Oxydome Contributors}
}
\date{\today}

\theoremstyle{definition}
\newtheorem{definition}{Definition}
\newtheorem{invariant}{Invariant}

\newcommand{\ox}{\textsf{Ox}}
\newcommand{\gas}{\textsf{GAS}}
\newcommand{\ohm}{\textsf{OH\scriptsize M}}\newcommand{\evk}{\textsf{EVK}}
\newcommand{\evfs}{\textsf{EVFS}}
\newcommand{\ovm}{\textsf{OVM}}
\newcommand{\ohms}{OHMs}
\newcommand{\bft}{\textsf{BFT}}
\newcommand{\poa}{\textsf{PoA}}
\newcommand{\outpoint}{\textsf{OutPoint}}
\newcommand{\pk}{\textsf{pk}}
\newcommand{\txid}{\textsf{txid}}
\newcommand{\WithdrawRequest}{\textsf{WithdrawRequest}}
\newcommand{\ExpressBusWithdraw}{\textsf{ExpressBusWithdraw}}
\newcommand{\BusReceipt}{\textsf{BusReceipt}}
\newcommand{\ChainId}{\textsf{ChainId}}
\newcommand{\ChainParams}{\textsf{ChainParams}}
\newcommand{\ConsensusProof}{\textsf{ConsensusProof}}
\newcommand{\DAO}{DAO}
\newcommand{\DAOAction}{\textsf{DAOAction}}
\newcommand{\EpochTransition}{\textsf{EpochTransition}}
\newcommand{\FileEvkLabel}{\textsf{FileEvkLabel}}
\newcommand{\InboundReceipt}{\textsf{InboundReceipt}}
\newcommand{\LockedReward}{\textsf{LockedReward}}
\newcommand{\Message}{\textsf{Message}}
\newcommand{\OhmState}{\textsf{OhmState}}
\newcommand{\OutboundMessage}{\textsf{OutboundMessage}}
\newcommand{\PriceAsset}{\textsf{PriceAsset}}
\newcommand{\PriceBounds}{\textsf{PriceBounds}}
\newcommand{\ProofManager}{\textsf{ProofManager}}
\newcommand{\SecureCommunicationLayer}{\textsf{SecureCommunicationLayer}}
\newcommand{\StorageLabel}{\textsf{StorageLabel}}
\newcommand{\VerificationCoordinator}{\textsf{VerificationCoordinator}}
\newcommand{\tx}{tx}
\newcommand{\hash}{\mathcal{H}}
\newcommand{\hashvalue}{\mathsf{Hash}}
\newcommand{\sign}{\mathsf{Sign}}
\newcommand{\verify}{\mathsf{Verify}}
\newcommand{\emptyroot}{\mathbf{0}}
\newcommand{\FP}{10^{12}}
\newcommand{\bps}{\mathrm{bps}}

\begin{document}

\maketitle

\begin{abstract}
Oxydome (the \ox{} network) is a decentralized ledger that combines a permissioned
Proof-of-Authority consensus core with an authenticated sparse-Merkle UTXO
state, a deterministic bytecode virtual machine, tenant-owned private chains,
and a family of fuel currencies.  This paper describes the protocol-level
systems of the Oxydome blockchain as they are realized in its reference
implementation: the cryptographic primitives, the authenticated state model, the
block and transaction format, the value-index currency engine, the layered
consensus and master-audit architecture, the \ohm{} tenant-chain subsystem, the
GAS and Speedometer fuel currencies, a coin-toss transaction-ordering
layer that deterministically labels every transaction for OVM execution,
the DAO governance framework, pegged
currencies and liquidity pools, cross-chain bridging via the Express Bus and
bridge module, and the ephemeral-verification-key infrastructure for node
authentication and secure channels.  All on-chain arithmetic is integer
fixed-point with no floating point, and every computation is deterministic so
that all honest nodes derive byte-identical results.
\end{abstract}

\tableofcontents
\newpage

\section{Introduction}

The \ox{} network (``Oxydome''), whose native unit is the \ox{} token (symbol
\ox{}), is a layered blockchain protocol designed around the principle of
\emph{authenticated delegation}: a small set of \emph{master} nodes does not
store block bodies or the UTXO set, yet can cryptographically audit every block
through Merkle inclusion proofs, a deterministic state-transition replay, and a
supply-conservation invariant.  Server nodes store and execute the full state,
produce blocks, and participate in a Practical-Byzantine-Fault-Tolerance-style
(PBFT-inspired) voting protocol layered on top of a round-robin leader
selection.

In addition to the base ledger, the protocol hosts several integrated subsystems:

\begin{itemize}[leftmargin=1.5em]
  \item The \textbf{Ox Virtual Machine} (\ovm{}): a deterministic, fixed-width
    bytecode interpreter for on-chain programs, supporting two deployment
    tiers (raw bytecode and compiled Solid-lang contracts).
  \item The \textbf{GAS} fuel currency: an account-based token minted from server
    rewards and airdropped according to an on-chain price oracle.
  \item The \textbf{Speedometer} stabilization currency: an account-based
    token minted per-transaction and airdropped proportionally to balances.
  \item \textbf{OHMs} (Ox Hosted Machines): tenant-owned private chains that
    execute privately but anchor their continuity into \ox{} blocks.
  \item A \textbf{DAO} framework: on-chain governance with maker-controlled
    currencies and pegged redemptions.
  \item \textbf{Express Bus} and \textbf{Bridge} modules: tracked
    cross-system value flow and verifiable cross-chain asset transfer.
  \item The \evk{} system: ephemeral per-block verification keys for node
    authentication, secure channels, and file-system labeling.
\end{itemize}

This paper is organized as follows.  Sections~\ref{sec:prelim}--\ref{sec:net}
cover the foundational systems; Section~\ref{sec:ovm} describes the virtual
machine; Section~\ref{sec:ohm} the tenant-chain system; and Section~\ref{sec:analysis}
discusses security and invariants.  Section~\ref{sec:coin-toss} documents
the coin-toss transaction-ordering layer.

\section{Prerequisites}
\label{sec:prelim}

\subsection{Cryptographic Primitives}
\label{sec:crypto}

All hashing in \ox{} uses \textbf{SHA3-512} (FIPS~202 / Keccak), producing
64-byte digests.  The protocol defines three derived operations:

\paragraph{Plain hash.} $\hash(x) = \mathrm{SHA3\text{-}512}(x)$, mapping an
arbitrary byte string to a $64$-byte \emph{hash}.

\paragraph{Length-prefixed multi-part hash.}
$\hash_{\mathrm{many}}(p_1,\dots,p_n)$ concatenates each part preceded by its
length in little-endian bytes, then hashes the result.  This prevents
concatenation ambiguity (``$ab$'' vs.\ ``$a$, $b$''):
\[
\hash_{\mathrm{many}}(p_1,\dots,p_n)
  = \mathrm{SHA3\text{-}512}\!\Bigl(\sum_{i=1}^{n}
  \bigl(\mathrm{LE64}(|p_i|)\,\|\,p_i\bigr)\Bigr).
\]

\paragraph{Serialized-value hash.}
$\hashvalue(v) = \hash(\mathsf{encode}(v))$, where $\mathsf{encode}(v)$ is a
deterministic serialization.  This is the primitive used to commit structured
values (leaves, proposals, signing payloads) to a hash.

\paragraph{Signatures.}
All signing and verification uses \textbf{Ed25519}.  A key pair consists of a
32-byte public key and a 64-byte signature.  $\verify(pk, m, \sigma)$ returns
true iff $\sigma$ is a valid Ed25519 signature on message $m$ under public key
$pk$.  The same primitive signs transactions, blocks, admission certificates,
governance votes, and audit attestations.

\paragraph{Position-dependent labeling hash.}
For transaction labeling (Section~\ref{sec:coin-toss}) the protocol also uses
\textbf{BLAKE2b-512}---the same 64-byte digest width as SHA3-512 but a
cryptographically secure, deliberately less common alternative to the SHA
family.  The \emph{position label} for transaction $j$ in block $h$ is
\[
\lambda_j = \mathrm{BLAKE2b\text{-}512}\!\Bigl(
  \tau,\; \mathrm{LE64}(j),\; \mathrm{LE64}(\tau_{\mathrm{block}}),\;
  \mathrm{LE64}(|\mathcal{M}|)\,\|\,\mathcal{M}
\Bigr),
\]
where $\tau = \hash(H_h)$ is the block header hash, $\tau_{\mathrm{block}}$ is
the block timestamp, and $\mathcal{M}$ is the network magic phrase.

\subsection{Merkle Trees}
\label{sec:merkle}

Two Merkle-tree constructions are used.

\paragraph{Standard Merkle tree.} Used for transaction roots, operation roots,
OVM-receipt roots, and the masters-root commitment.  Leaves are
$\mathbf{0x00}\,\hash(\ell)$ and internal nodes are
$\mathbf{0x01}\,\hash(\ell)\,\hash(r)$, giving second-preimage resistance.
An odd level duplicates its last node.  The empty set commits to
$\emptyroot$ (the $64$-byte all-zero hash).  An inclusion proof is the list of
siblings from the leaf level to the root; verification recomputes the path and
checks equality with the stored root.

\paragraph{Sparse Merkle tree.}
Used for the UTXO ledger.  It is a $256$-level binary tree
($D = 256$) in which every leaf sits at exactly depth $D$.  An empty subtree
commits to $\emptyroot$ at every level.  Keys are derived as
\[
\mathsf{path}(o) = \mathrm{first\_32\_bytes}\bigl(\hashvalue(o)\bigr),
\]
where $o$ is an \outpoint\ (Definition~\ref{def:outpoint}).  A leaf stores a
\emph{leaf commitment} $\hashvalue(\mathit{entry})$, where an entry is
$(\mathit{outpoint}, \mathit{output}, \mathit{created\_height})$.  Inclusion
proofs have exactly $D$ siblings ordered leaf-to-root; verification is purely a
function of the sibling list and the path bits.  The tree is \emph{pure}: the
root depends only on the set of live entries, not their insertion order, which
makes it portable between implementations.

\section{State Model}
\label{sec:state}

\subsection{The UTXO Set}

The primary on-chain state is an authenticated UTXO set implemented as a sparse
Merkle tree of depth $256$.  Each unspent output is an entry
\[
e = (o,\; v,\; h)
\]
where $o$ is an \outpoint, $v$ is a TxOutput, and $h$ is the height at which the
output was created.  The \outpoint\ (Definition~\ref{def:outpoint}) is a pair
$(t, i)$: a $64$-byte transaction identifier and a $32$-bit output index.

\begin{definition}[Outpoint]
\label{def:outpoint}
An \outpoint\ $o = (t, i)$ references output $i$ of the transaction whose
identifier is $t$.
\end{definition}

\paragraph{TxOutput.} A TxOutput is a tuple $(\mathit{amount}, \mathit{recipient}, \mathit{kind}, \mathit{rebound\_gen})$
where $\mathit{amount}$ is a non-negative integer, $\mathit{recipient}$ is a
32-byte public key (address), and $\mathit{kind} \in K$ (Section~\ref{sec:output-kinds}).

\subsection{Output Kinds}
\label{sec:output-kinds}

The \emph{kind} of an output determines how and by whom it may be spent:

\begin{itemize}[leftmargin=2em]
  \item \textbf{Spendable} --- ordinary transferable value.
  \item \textbf{LockedReward} --- node or master fee earnings; spendable only
    via an approved \WithdrawRequest.
  \item \textbf{ExitReserve($g$)} --- value deposited into the exit/reserve sink
    in generation $g$; exerts downward pressure on the value index.  Refunded
    (tagged) on a rebound.
  \item \textbf{DaoTreasury($d$)} --- held by a DAO until a DAO-approved
    transfer releases it.
  \item \textbf{OhmTreasury($m$)} --- held by an \ohm{}'s treasury; spendable
    only via OHM operations.
  \item \textbf{BusReceipt($\beta$, $a$)} --- tagged by an Express-Bus deposit;
    spendable only via an \ExpressBusWithdraw.
\end{itemize}

\paragraph{Rebound credits.}
Outputs created by a rebound event carry $\mathit{rebound\_gen} = g$, marking
them as maturity-locked.  A rebound-credited output created at height $h$ in
generation $g$ may not be spent before
\[
h + L \quad\text{or before generation } g{+}1,
\]
where $L$ is the rebound lock blocks parameter.  This ``bunny-hop'' maturity
prevents rapid respending of rebound refunds and stabilizes the value index.

\subsection{Composite State Root}
\label{sec:state-root}

The block header commits a single \emph{state root} that binds together the
UTXO tree and all optional subsystem states:
\[
R_{\mathrm{state}} = \hash\bigl(
  \mathit{domain},\;
  R_{\mathrm{utxo}},\;
  R_{\mathrm{ovm}},\;
  R_{\mathrm{dao}},\;
  R_{\mathrm{ohm}},\;
  R_{\mathrm{storage}},\;
  R_{\mathrm{timer}},\;
  R_{\mathrm{pegged}},\;
  R_{\mathrm{pool}},\;
  R_{\mathrm{file}},\;
  R_{\mathrm{express}}
\bigr).
\]
The domain-separator string (``\texttt{oxydome-combined-state-v9}'' family) is
chosen feature-sensitively so that chains running a subset of features retain
byte-identical roots to prior versions.  The header also commits the bare
UTXO root $R_{\mathrm{utxo}}$ to a separate field so that header-only masters
can verify UTXO-spend inclusion proofs without the composite root (which also
covers contract code and storage they do not store).

\section{Block Structure}
\label{sec:block}

A block $B$ at height $h$ consists of a header $H$ and a body
$(\mathcal{T}, \mathcal{O}, \mathcal{R})$:

\paragraph{Header fields.}
\[
H = \bigl(
  h,\; \nu,\; H_{\mathrm{ovm\_act}},\; H_{\mathrm{sc\_act}},\;
  H_{\mathrm{prev}},\; R_T,\; R_O,\; R_{\mathcal{R}},\;
  R_{\mathrm{utxo}},\; R_{\mathrm{state}},\;
  \Sigma,\; e,\; g,\; V,\; p,\; \tau,\;
  \Phi,\; n_T,\; \Phi_{\mathrm{exit}},\;
  s_{\mathrm{toss}},\; R_{\mathcal{T}}^{\mathrm{ord}}
\bigr).
\]
Here:
\begin{itemize}[leftmargin=2em]
  \item $\nu$ is the execution version ($1$ = legacy, $2$ = \ovm{}, $4$ =
    smart-contract artifact);
  \item $H_{\mathrm{ovm\_act}}$, $H_{\mathrm{sc\_act}}$ are activation heights;
  \item $H_{\mathrm{prev}} = \hash(H_{h-1})$ links blocks;
  \item $R_T, R_O, R_{\mathcal{R}}$ are Merkle roots over the transaction,
    operation, and OVM-receipt lists;
  \item $R_{\mathrm{utxo}}$ and $R_{\mathrm{state}}$ are the UTXO and composite
    state roots (Section~\ref{sec:state-root});
  \item $\Sigma$ is the circulating supply;
  \item $e$ is the epoch; $g$ is the rebound generation;
  \item $V$ is the value index (fixed-point, Section~\ref{sec:value-index});
  \item $p$ is the block producer (a server-node public key);
  \item $\tau$ is the timestamp;
  \item $\Phi$ is the total transaction volume (window statistic);
  \item $n_T$ is the transaction count; $\Phi_{\mathrm{exit}}$ is the exit sink amount;
  \item $s_{\mathrm{toss}}$ is the coin-toss seed commitment (Section~\ref{sec:coin-toss});
  \item $R_{\mathcal{T}}^{\mathrm{ord}}$ is the root of the ordered-transaction
    Merkle tree (Section~\ref{sec:coin-toss}).
  \item optional EVK extension fields (Section~\ref{sec:evk}).
\end{itemize}

\paragraph{Header roots must match.}
The producer computes $R_T$, $R_O$, $R_{\mathcal{R}}$ as Merkle roots over the
body, and $R_{\mathrm{utxo}}$, $R_{\mathrm{state}}$ from the resulting state.
Any node re-applying the block independently verifies
\[
R_T = \mathsf{txroot}(\mathcal{T}),\quad
R_O = \mathsf{oproot}(\mathcal{O}),\quad
R_{\mathcal{R}} = \mathsf{receiptroot}(\mathcal{R}),
\]
and, after coin-toss ordering (Section~\ref{sec:coin-toss}),
\[
R_{\mathcal{T}}^{\mathrm{ord}}
   = \mathsf{ordroot}(\mathcal{T}_{\pi}),
\]
where $\mathcal{T}_{\pi}$ is the producer-committed permutation of $\mathcal{T}$.

\subsection{State Transition}

Applying a block is a pure function
\[
\Sigma' = \mathsf{Apply}(\Sigma,\; B,\; G,\; P),
\]
where $G$ is the GovernanceSet and $P$ the ChainParams.  The transition:
\begin{enumerate}[label=(\arabic*),leftmargin=2em]
  \item validates linkage ($H_{\mathrm{prev}}$, height, execution version,
    producer authorization);
  \item consumes inputs against the UTXO set (verifying signatures and
    maturity for rebound credits);
  \item creates outputs (updating the UTXO tree, tracking rebound generations);
  \item executes governance operations (mint, burn, withdraw, rewards, DAO,
    OHM, GAS, Express Bus, liquidity pool, timer, file-system, speedometer,
    node-leave, departure-earnings);
  \item executes OVM receipts committed in $\mathcal{R}$;
  \item recomputes the value index, generation, and supply from window
    statistics and explicit operation amounts;
  \item commits the new $R_{\mathrm{utxo}}$ and $R_{\mathrm{state}}$.
\end{enumerate}
The \emph{lineage invariant} $\sum_{\mathit{live}} \mathit{amount} = \Sigma$
must hold after every transition (Section~\ref{sec:audit}).

\section{Transactions}
\label{sec:tx}

A transaction is a value transfer with optional program execution:
\[
\mathcal{T} = (I,\; O,\; \phi,\; \alpha),
\]
where $I$ is a list of inputs, $O$ a list of outputs, $\phi$ the fee, and
$\alpha$ an optional OVM action.

\paragraph{Input.} Each input $I_j = (o_j,\, pk_j,\, \sigma_j)$ references an
\outpoint\ $o_j$, declares its owner $pk_j$, and carries an Ed25519 signature
$\sigma_j$ over a canonical signing payload that excludes all signatures
(self-reference-free commitment).

\paragraph{Conservation.}
A transaction is valid iff the owner of every input holds a live UTXO at that
\outpoint, every input signature verifies, no input is spent twice, and
\[
\sum_j \mathit{amount}(o_j) \;=\; \sum_i \mathit{amount}(O_i) \;+\; \phi.
\]

\paragraph{Fee.}
The base transfer fee is
\[
\phi_{\mathrm{base}} = \max\!\Bigl(1,\; \bigl\lfloor
  \textstyle\sum_i \mathit{amount}(O_i) \cdot 12 / 10^{6}\bigr\rfloor\Bigr),
\]
scaled by the current value index at block construction time.

\paragraph{OVM action.}
An action $\alpha$ attached to a transaction requests on-chain program execution:
\begin{itemize}[leftmargin=2em]
  \item \textbf{Deploy}: upload raw bytecode (a $\leq 64\,\mathrm{KiB}$ program)
    and reserve $\leq 2\,\mathrm{GiB}$ of data stack.
  \item \textbf{Call}: invoke a deployed contract by id.
  \item \textbf{DeployContract}: deploy a compiled Solid-lang contract artifact
    ($\leq 256\,\mathrm{KiB}$).
  \item \textbf{CallContract}: invoke a named entry point on a contract, with
    an authority-input word and up to 64 arguments.
\end{itemize}
Gas for the action is bounded by $\lceil\mathsf{stack\_bytes}/6\rceil$ per
transaction and by an aggregate block limit.

\paragraph{Position hash.}
Once a block producer has committed to an ordering of $\mathcal{T}$, each
transaction is labeled with a deterministic \emph{position hash}:
\[
\hash_{\mathcal{T}} = \mathrm{BLAKE2b\text{-}512}\!\Bigl(
  \hash(H),\; \mathrm{LE64}(j),\; \mathrm{LE64}(\tau),\; \mathcal{M}
\Bigr),
\]
where $j$ is the transaction's position in the ordered list, $\hash(H)$ is
the block header hash, $\tau$ is the block timestamp, and $\mathcal{M}$ is the
network magic phrase.  This value is injected into the OVM execution context so
that on-chain programs can derive a unique, bias-resistant, reproducible
position label without trusting the producer.

\section{Value Index and Currency Mechanics}
\label{sec:value-index}

\subsection{Fixed-Point Arithmetic}

All currency computation uses signed $128$-bit fixed point with scale
$\FP = 10^{12}$, so $\FP$ represents $1.000$.  Multiplication and division are
saturating; a deterministic natural-logarithm (\texttt{ln\_fixed}) and
exponential (\texttt{exp\_fixed}) are provided via range reduction and the
\href{https://en.wikipedia.org/wiki/Rapidly_convergent_statements_for_ln}{atanh series}
\[
\ln(m) = 2\sum_{k=0}^{\infty} \frac{z^{2k+1}}{2k+1},\qquad z=\frac{m-1}{m+1},\quad m\in[1,2).
\]
No floating-point arithmetic is used on-chain, guaranteeing determinism across
platforms.

\subsection{Window Statistics}

Each block records $\omega_h = (\Phi_h,\, n_{T,h},\, \Phi_{\mathrm{exit},h},\,
\Phi_{\max,h},\, f_h)$ where $\Phi$ is volume, $n_T$ the transaction count,
$\Phi_{\mathrm{exit}}$ the exit-sink amount, $\Phi_{\max}$ the highest
transaction output total, and $f_h$ the Express-Bus flow signal.  The rolling
window $W = \sum_{j=h-W_0+1}^{h} \omega_j$ aggregates these over the last
$P = \texttt{block\_window}$ blocks.

\subsection{Index Formula}
\label{sec:value-formula}

The value index $V_h$ updates via an exponential moving average:
\[
V'_h = \lambda\cdot v_0 \cdot \frac{A_h}{X_h} + (1-\lambda)\cdot V_{h-1},
\]
where $\lambda = \lambda_{\bps}/10^{4}$ (default $0.10$), $v_0$ is the initial
index ($\FP$ at genesis), and

\begin{align}
A_h &= 1 + \alpha\,\ln\!\Bigl(1+\tfrac{T}{t_{\mathrm{ref}}}\Bigr)
        + \beta\,\ln\!\Bigl(1+\tfrac{A_{\mathrm{mean}}}{a_{\mathrm{ref}}}\Bigr),
\\
X_h &= \gamma\cdot\frac{E}{e_{\mathrm{ref}}}\cdot\Bigl(\frac{S_0}{S_h}\Bigr)^{\eta},
\\
f_h &= \xi\cdot f_{\mathrm{express},h}.
\end{align}

Here $T = \Phi_W$ is window volume, $A_{\mathrm{mean}} = T/n_T$, $E=\Phi_{\mathrm{exit},W}$,
$S_0$ is genesis supply, $S_h$ is current supply, and $f_{\mathrm{express}}$
is the Express-Bus net-flow signal.  The weights $\alpha,\beta,\gamma,\eta,\xi$
are basis-point parameters in $\ChainParams$.

\paragraph{Rebound.}
If $V'_h \le 0$, the chain \emph{rebounds}: it resets
$V_h = V_{h-1}$ (the last positive value) and increments the generation $g$.
Rebound credits created at generation $g$ are maturity-locked until
generation $g+1$ (Section~\ref{sec:bunnyhop}).

\subsection{Guardrails}
\label{sec:guardrails}

Governance actions are constrained to keep the value index stable:
\begin{itemize}[leftmargin=2em]
  \item \textbf{Value-impact limit}: any single action may move $V$ by at most
    $\delta = \texttt{value\_impact\_limit\_bps}/10^{4}$ of its prior value.
  \item \textbf{Supply-change cap}: mint/burn per epoch $\le$
    $\epsilon\cdot S_0$ where $\epsilon = \texttt{supply\_change\_limit\_bps}/10^{4}$
    (default $1\%$).
  \item \textbf{Fee affordability}: $\phi_{\mathrm{eff}} \le \rho\cdot A_{\min}$.
  \item \textbf{Withdrawal cap}: server withdrawals per epoch $\le \kappa\cdot S_0$.
  \item \textbf{Per-account}: at most $\pi = 30\%$ of an account's balance
    withdrawable per operation.
\end{itemize}

\section{Consensus}
\label{sec:consensus}

\ox{} combines Proof-of-Authority (\poa{}) leader selection with a two-phase
BFT voting protocol.

\subsection{Leader Selection}

For view $v$ at height $h$, the leader is
\[
\mathsf{leader}(v) = \mathrm{nodes}[v \bmod |\mathrm{nodes}|],
\]
taken from the current server-node set.  Only the leader may propose a block at
its view.

\subsection{BFT Voting}

Two vote kinds are defined:

\begin{description}
  \item[Prep\text{{}are}] cast once a node has validated the proposal and
    replayed the state transition;
  \item[Comm\text{}it] cast once the Prepare quorum certificate (2/3 super-majority)
    is observed.
\end{description}

Each vote $\upsilon = (h, v, B, \kappa, \pi, \sigma_\pi)$ signs
$\hash(h,\,v,\,\hash(B),\,\kappa)$ under voter key $\pi$, where $\kappa$ is the
phase.  A \emph{quorum certificate} $Q_\kappa$ collects votes for a common
$(h,v,B,\kappa)$ and is valid when the count of distinct, authorized, validly
signed voters is $\ge \texttt{quorum}(n)$, where
\[
\texttt{quorum}(n) = \frac{n\cdot\texttt{num}}{\texttt{den}} + 1.
\]
With the default $\frac{2}{3}$ threshold and $n$ nodes, a quorum is
$\lfloor 2n/3\rfloor + 1$.

\subsection{Master Check}

At each epoch boundary, server nodes issue a \emph{MasterCheck} attestation
$\mu = (e,\, R_M,\, \pi,\, \sigma_\pi)$ signing the epoch number $e$ and the
\emph{masters root}
\[
R_M = \mathsf{MerkleRoot}\bigl(\{\;\hash(\texttt{``master''}\,\|\,\pk)\;|\;\pk\in\mathrm{masters}\;\}\bigr),
\]
with leaves sorted by key.  The epoch is accepted only when
$\texttt{quorum}(n)$ distinct nodes attest the same $(e, R_M)$.

\subsection{Epoch Governance}
\label{sec:epoch-gov}

Master-set changes proceed in two stages:
\begin{enumerate}[label=(\arabic*),leftmargin=2em]
  \item \textbf{Proposal}: a server node proposes a candidate master key.
    A strict majority of server nodes sign the proposal id.
  \item \textbf{Transition}: the outgoing master set co-signs an
    \EpochTransition, valid only with a $\frac{2}{3}$ (default) super-majority
    of outgoing master signatures.
\end{enumerate}
Both stages are broadcast over the wire (Section~\ref{sec:net}).

\section{Coin-Toss Transaction Ordering}
\label{sec:coin-toss}

Because the leader is a single server node, the \emph{order} in which it
places transactions in a block $\mathcal{T}$ is a vector of trust.  A
malicious or biased leader could reorder transactions to front-run,
censor, or sandwich other participants' transfers.  \ox{} mitigates this
risk with a \emph{coin-toss} ordering layer that runs \emph{in parallel} to
the BFT consensus flow: every honest node independently derives the same
permutation of $\mathcal{T}$, so the producer cannot choose which permutation
its block commits to.

\subsection{Seed Derivation}
\label{sec:coin-toss-seed}

The coin-toss seed for block $h$ is an epoch-global random value committed
to on-chain at the end of the previous epoch and revealed at epoch
transition time.  Concretely, during epoch $e-1$ each master $\pi_i$
contributes a commitment
\[
c_i = \hash(\texttt{``toss''}\,\|\,\pk_i\,\|\,r_i),
\]
where $r_i$ is a 256-bit secret.  The commitments are aggregated into a
Merkle root $R_{\mathrm{seeds}}$ stored in the final block header of the
epoch.  At the epoch-$e$ transition, each master reveals $r_i$; the seed
for epoch $e$ is then
\[
s_e = \hash\!\Bigl(\sum_{i} \bigl(\mathrm{LE64}(|r_i|)\,\|\,r_i\bigr)\Bigr).
\]
The header field $s_{\mathrm{toss}}$ stores $\hash(s_e \,\|\, \tau)$---a
commitment binding the seed to the block timestamp so that the producer
cannot retroactively adjust $\tau$ to influence the draw.

\subsection{Permutation Algorithm}
\label{sec:coin-toss-perm}

Given the seed-derived digest $\tau = \hash(s_{\mathrm{toss}})$ (a 64-byte
value) and a block with $n$ transactions, the canonical permutation $\pi$
of $\{0,\dots,n-1\}$ is computed as a \emph{Fisher--Yates shuffle} driven by
an AES-256-CTR stream cipher seeded with $\tau$:

\begin{enumerate}[label=(\arabic*),leftmargin=2em]
  \item Interpret $\tau$ as a 512-bit seed.  Derive the 256-bit AES key
    $k = \tau[0:32]$ and the 128-bit initial counter block $c_0 = \tau[32:48]$.
  \item Generate $\lceil n\log_2 n\rceil$ bytes of keystream by encrypting
    successive counter values $c_0, c_0{+}1, c_0{+}2,\ldots$ under AES-256-CTR
    with key $k$.
  \item For $i$ from $n-1$ down to $1$, draw index $j$ uniformly from
    $\{0,\dots,i\}$ using the next $\lceil\log_2(i+1)\rceil$ unbiased bits
    from the keystream (rejection-sampled if the drawn value exceeds $i$),
    and swap positions $i$ and $j$.
\end{enumerate}

The result is a bijection $\pi : \{0,\dots,n-1\} \to \{0,\dots,n-1\}$ that
every node computes identically.  The block producer commits to
$R_{\mathcal{T}}^{\mathrm{ord}} = \mathsf{ordroot}(\mathcal{T}_\pi)$ in the
header, where $\mathcal{T}_\pi$ denotes the transaction list rearranged
according to $\pi$.  Any node that re-derives $\pi$ from the published
seed and re-applies the shuffle obtains the same root and can reject blocks
that deviate.

\subsection{Non-Interactive and Bias-Resistant}

The protocol is \emph{non-interactive}: no node needs to send anything
during block construction to influence the draw.  The seed comes from the
previous epoch's revealed master contributions, which are committed before
the block in question is ever proposed.  Because at least one master is
assumed honest (Section~\ref{sec:audit}), the seed is unbiased with
overwhelming probability, and a malicious leader cannot adapt the seed to
favor a particular ordering.

\subsection{Relation to Consensus}

The coin-toss layer does not replace BFT consensus; it constrains the
\emph{proposal} phase.  The leader still proposes the unordered set
$\mathcal{T}$ and the execution result, but the ordering it commits to in
$R_{\mathcal{T}}^{\mathrm{ord}}$ must match the deterministic coin-toss
permutation.  Validators check this root during Prepare voting (Section~\ref{sec:consensus}) and include its verification in the block state-transition
checks (Section~\ref{sec:block}).  If a leader submits a
mismatched root, the block is rejected outright.

\section{Node Admission and Software Integrity}
\label{sec:admission}

\subsection{Code Manifest}

A node's software is summarized by a \emph{CodeManifest}: a sorted map from file
path to $\hash(\text{file bytes})$.  A single manifest root
\[
R_{\mathrm{manifest}} = \mathsf{MerkleRoot}\bigl(\{\;\hash(\mathit{path}\,\|\,\hash(\mathit{file}))\;\}\bigr)
\]
commit to the entire on-disk program so that any unmodified-binary proof
reduces to a Merkle inclusion proof against this root.

\subsection{Admission Certificate}

A server node joins an epoch only if a master has issued an
\emph{AdmissionCertificate}
\[
c = (\pi_{\mathrm{node}},\, R_{\mathrm{manifest}},\, e,\, t_{\mathrm{issue}},\, t_{\mathrm{expiry}},\, \pi_{\mathrm{master}},\, \sigma_{\pi_{\mathrm{master}}}),
\]
where the master $\pi_{\mathrm{master}}$ signs the commitment
$\hash(c_{\mathrm{signed}})$.  The certificate is valid iff the master is a
current epoch master, the manifest root is in the approved set, and
$t_{\mathrm{expiry}}$ has not passed.  Masters maintain the approved-manifest
set and a revocation list.

\section{Audit Architecture}
\label{sec:audit}

\ox{} separates the full-validation (``server'') role from the
header-only (``master'') role through a layered audit.

\subsection{Per-Block Master Audit}

For each proposed block, the proposer supplies an \emph{AuditPackage}
\[
\mathcal{P} = (B,\; R_{\mathrm{utxo},h{-}1},\; \{\,\pi_{\mathrm{in}}\,))
\]
containing Merkle inclusion proofs $\pi_{\mathrm{in}}$ for every transaction
input against the prior UTXO root.  A master then:
\begin{enumerate}[label=(\arabic*),leftmargin=2em]
  \item verifies every $\pi_{\mathrm{in}}$ against $R_{\mathrm{utxo},h-1}$;
  \item checks the proven entry matches the spent \outpoint{} and recipient;
  \item checks rebound maturity of inputs;
  \item replays the state transition and checks \textbf{conservation} per
    transaction (${\sum I = \sum O + \phi}$);
  \item verifies header statistics (volume, count, exit amount) match the body;
  \item verifies total gas fees $\sum r_{\mathrm{gas}} \cdot p_{\mathrm{gas}}$;
  \item recomputes supply, value index, and generation from operations and
    checks them against the header;
  \item enforces all guardrails (supply cap, value-impact, withdraw caps).
\end{enumerate}
The master signs the resulting \emph{AuditReport}
$(\hash(B),\, R_{\mathrm{state}},\, R_{\mathrm{audit}})$, where
$R_{\mathrm{audit}}$ is the Merkle root over sorted input-leaf hash commitments.
A block is accepted when $\ge\texttt{quorum}(|\mathrm{masters}|)$ audits are
collected (\textbf{attestation quorum}).

\subsection{Lineage Invariant}
\begin{invariant}[Lineage]
At every block boundary,
\[
\sum_{\mathit{live\;UTXO}} \mathit{amount} \;=\; \Sigma_h.
\]
\end{invariant}
This is checked by both server replay and master audit.  A violation
($\sum\neq\Sigma$) is a fatal state error.

\subsection{Deep Replay}

Every \texttt{audit\_checkpoint\_blocks} ($K$, default $50$) blocks, server
nodes holding full state perform a \emph{literal deep replay}: they re-apply
$K$ blocks from scratch and compare the resulting $R_{\mathrm{utxo}}$ to the
committed value.  This is the only check that independently recomputes the UTXO
tree root (the per-block audit trusts the proposer's root but verifies
supply/index/statistics).

\section{The Ox Virtual Machine}
\label{sec:ovm}

\subsection{Architecture}

The \ovm{} is a deterministic, stack-based bytecode interpreter.  A program
consists of a fixed-width operand stack, indexed local variables, $u_64$ storage
slots, and a bounded data stack.  All arithmetic is on $u_{64}$ words with no
overflow widening, and the interpreter enforces hard caps on stack depth
($1\,024$), locals ($256$), storage slots ($4\,096$), and contracts
($10\,000$).  The instruction set is dispatched by an inline-interpreter loop;
arithmetic opcodes compile to native machine instructions at build time while
remaining bit-identical across conformant implementations.

\subsection{Instructions}

Instructions operate on typed $u_{64}$ values and include:
\begin{itemize}[leftmargin=2em]
  \item \texttt{Push}, \texttt{LoadArg}, \texttt{LoadLocal}, \texttt{StoreLocal},
    \texttt{LoadStorage}, \texttt{StoreStorage}, \texttt{LoadStorageDyn},
    \texttt{StoreStorageDyn};
  \item \texttt{Add}, \texttt{Sub}, \texttt{Mul}, \texttt{Div}, \texttt{Mod},
    \texttt{Equal}, \texttt{LessThan}, \texttt{GreaterThan}, \texttt{And},
    \texttt{Or}, \texttt{Not}, \texttt{BitXor}, \texttt{Shl}, \texttt{Shr};
  \item \texttt{Jump}, \texttt{JumpIfFalse}, \texttt{Emit};
  \item \texttt{Return}, \texttt{Revert};
  \item \texttt{XCall}, \texttt{ReadReply} --- contract-originated cross-\ohm{} messaging;
  \item \texttt{LoadPriceBounds}, \texttt{PeggedConvert},
    \texttt{LoadPoolState}, \texttt{PoolSwap} --- host oracles;
  \item \texttt{DataSize}, \texttt{LoadData}, \texttt{StoreData} --- data stack I/O.
\end{itemize}

The \texttt{Hash} instruction consumes $k$ words (top-first), takes
$\mathrm{SHA3\text{-}512}$ of their little-endian byte concatenation, and
pushes the first $8$ bytes as a $u_{64}$ selector (used for mapping-slot
derivation).

\subsection{Gas}

Each OVM action reserves
\[
g = \lceil \mathsf{stack\_bytes} / 6\rceil
\]
units of gas, with a per-transaction cap and a block-level aggregate cap
($10^{7}$ units).  Execution charges gas equal to instructions used; failures
(OOM, division by zero, overflow, revert, out-of-gas) produce a stable,
consensus-visible failure code so auditors need not re-execute bytecode to
confirm a failure.

\subsection{Contract Deployment and Storage}

Two deployment tiers exist:
\begin{description}
  \item[Legacy (v2, \ovm)] raw bytecode $\leftrightarrow$
    $\mathrm{id} = \hash(\texttt{``oxydome-contract-v1''}\,\|\,\txid)$.
  \item[Smart-contract (v4)] compiled \textbf{Solid-lang} contract artifacts,
    each containing named entry points.  Entry points are invoked via
    \texttt{CallContract}, which receives authority words and arguments.
\end{description}
A contract has shared $u_{64}\!\to\!u_{64}$ storage ($4\,096$ slots max).  All
contract state lives under the OVM-state component of $R_{\mathrm{state}}$
(Section~\ref{sec:state-root}).

\subsection{Solid-lang}

The \emph{Solid} compiler front-end parses a practical subset of Solidity
syntax (contracts, interfaces, libraries, inheritance, mappings, structs,
enums, modifiers, events, custom errors, control flow) and compiles a
restricted native subset to \ovm{} bytecode: $u_{64}$ and \texttt{bool} scalars,
scalar state, supported arithmetic/logical expressions, assignments,
conditionals, and loops.  It additionally compiles \DAO{} policies, liquidity-pool
configurations, and pegged-currency configs into canonical on-chain encodings.

\subsection{Host Oracles}

The \ovm{} execution environment exposes optional \emph{host oracles} that
contracts may consult:
\begin{description}
  \item[PriceOracle] maps an asset id to packed \PriceBounds, used by
    \texttt{LoadPriceBounds} and \texttt{PeggedConvert};
  \item[PoolOracle] maps a pool id to packed reserves, $K$, and fee, used by
    \texttt{LoadPoolState} and \texttt{PoolSwap};
  \item[CrossHost] carries the originating \ohm{}, an inbox of cross-\ohm
    message receipts, and the oracles above, enabling \texttt{XCall} and
    \texttt{ReadReply}.
\end{description}
All oracle values are packed to $u_{64}$ via integer division by $10^6$ for
\ovm{} consumption.

\section{DAO Framework}
\label{sec:dao}

\subsection{DAO Policy}

A DAO is created by submitting a \DAOAction::Create with a canonical
\texttt{ODAO$\backslash$0x01}-prefixed policy encoding a sorted maker set
($\le 128$), an optional voting currency (1--12 uppercase ASCII characters,
positive initial supply), and an optional peg (target DAO, reduced-ratio
numerator/denominator).  The DAO id is
\[
d = \hash(\texttt{``oxydome-dao-id-v1''}\,\|\,\mathit{id}\,\|\,\mathit{policy}).
\]

\subsection{DAO State and Governance}

Each DAO maintains its own currency supply, balances, treasury, locked
reserves, pegged reserves, proposals, and a local value index.  Proposals
carry one of several actions (treasury transfer, mint, burn, set-peg) and
progress through holder-voting and maker-voting phases before execution or
expiration.  Two-tier voting requires both holder and maker majorities.

\subsection{Pegged Currencies}
\label{sec:pegged}

A pegged currency publishes a price band with upper bound, lower bound,
midpoint, and current price, recomputed as a moving average over
$\texttt{window}$ blocks with a $\texttt{spread}$ (default $5\%$).  A
\emph{guaranteed exchange rate} applies when the current price is in the upper
half of the band (${\ge \mathrm{midpoint}}$): the pegged token equals the main
currency at a $1{:}1$ rate.  When out of band, conversions fail.  These bounds
feed the GAS price multiplier (Section~\ref{sec:gas}).

\subsection{Liquidity Pools}
\label{sec:pools}

Pools implement a constant-product AMM $x \cdot y = K$.  Reserves and $K$ are
stored at $\FP$ scale; $K$ is recomputed each block in log-space as
\[
\ln K = 2\,V_{\mathrm{norm}}\,\ln(\text{upper}) +
        \ln(V_{\mathrm{norm}})\,\ln(\text{lower}),
\]
where $V_{\mathrm{norm}} = \mathrm{volume}/t_{\mathrm{ref}}$ is the
fixed-point ratio of accumulated trading volume to a reference.  The packed
swap executes the standard constant-product formula with a $\texttt{fee\_bps}$
deduction (default $100$~bps) on the input before ratio change:
\[
\Delta y = y - \frac{K}{x + \Delta x_{\mathrm{eff}}}.
\]

\section{GAS and Speedometer Fuel Currencies}
\label{sec:fuel}

\subsection{GAS}
\label{sec:gas}

\gas{} is the fuel token for on-chain program execution.  Key properties:
\begin{itemize}[leftmargin=2em]
  \item \textbf{Minting:} $20\%$ of every server-node reward is burned as \ox{}
    and minted as an equal amount of \gas{} to the same node.
  \item \textbf{Airdrop:} every \texttt{airdrop\_interval\_blocks} blocks, fresh
    \gas{} is distributed proportionally to balances, sized by
    \[
    \mathit{airdrop} = \Phi_{\max}\cdot\frac{\texttt{airdrop\_base\_bps}}{10^4}\cdot
      \frac{\pi_{\mathrm{ox}}}{\pi_{\mathrm{ref}}}.
    \]
  \item \textbf{Price oracle:} the \ox{} price $\pi_{\mathrm{ox}}$ is a
    master-governed, smoothed oracle bounded by
    $\texttt{price\_change\_limit\_bps}$ per update interval.
  \item \textbf{Burn:} \gas{} is consumed (burned) when a program is uploaded
    or executed.
\end{itemize}
\gas{} uses an account-balance ledger with per-account nonces (replay
protection) but is \emph{not} committed into $R_{\mathrm{state}}$.

\subsection{Speedometer}
\label{sec:speedometer}

The Speedometer currency is an optional stabilization layer:
\begin{itemize}[leftmargin=2em]
  \item \textbf{Per-tx mint:} a fixed share
    ($\texttt{speedometer\_stabilization\_bps} = 10$~bps) of each transaction's
    output total is freshly minted and accumulated.
  \item \textbf{Airdrop:} every
    $\texttt{speedometer\_airdrop\_interval\_blocks}$ blocks, minted amounts are
    distributed proportionally to current Speedometer balances, scaled by the
    Ox price oracle.
  \item \textbf{Accounts:} account-based, signed with nonces, not committed
    into $R_{\mathrm{state}}$ (like \gas{}).
\end{itemize}

\section{Ox Hosted Machines (Tenant Chains)}
\label{sec:ohm}

\subsection{Design}

An \ohm{} is a tenant-owned private blockchain whose transaction contents are
private yet whose \emph{continuity} is anchored into \ox{} blocks.  A tenant
builds private \ohm{} blocks and periodically submits an
\texttt{OhmAction::Update} operation that commits the new OHM context
(height, previous-hash, timestamp, state root) and the block hash to the \ox{}
ledger.  Server nodes execute the OHM block body off-radar, re-derive the
OHM state root, and reject any \ox{} block whose anchor does not match.

\subsection{OHM State}

The Ox chain stores, per \ohm{} $m$, only a lightweight \OhmState (policy, current
height, previous hash, anchored state root, supply, treasury balance).  The
live OHM UTXO set and contract state are re-derived by server nodes from the
private block bodies; only the state root is a consensus commitment.  The well-known
treasury outpoint $\mathsf{ohm\_treasury\_outpoint}(m) = (h, 0)$ holds the
\ox{} backing for the OHM treasury.

\subsection{Fee Flow}

Every \ohm{} transaction pays a fee that is split automatically:
\[
(\text{server share},\; \text{ohm share}) =
\textstyle\frac{40\%}{\text{of fees}}\;\;/\;\;\frac{60\%}{\text{of fees}}.
\]
The server share is distributed as \LockedReward{} outputs (participation-weighted
with a $3{:}1$ proposer bonus); the ohm share stays in the OHM treasury.

\subsection{Cross-OHM Messaging}
\label{sec:xoh}

\ohms{} communicate via a store-and-forward inbox/outbox:
\begin{itemize}[leftmargin=2em]
  \item \textbf{Outbound:} a contract stages an \texttt{XCall} in \ovm{}; the
    ledger enqueues an \OutboundMessage into the source OHM's outbox, committed
    to the OHM state root.  The message id is
    $\mathsf{msg\_id} = \hash(\texttt{``oxydome-xchain-msg-v1''}\,\|\,m\,\|\,h\,\|\,\mathrm{seq})$.
  \item \textbf{Relay:} the receiving OHM's maintainer attaches a
    \texttt{CrossRecvCall} to a transaction, causing the node to execute the
    target entry point and record an \InboundReceipt{} (with outcome and optional
    reply id).
  \item \textbf{Replies:} if the callee set \texttt{want\_reply}, an automatic
    reply message carrying $(\mathsf{msg\_id},\mathrm{success},\mathrm{value})$
    is relayed back.
  \item \textbf{Transfers:} a \texttt{CrossTransferCall} locks value on the
    source and mints a wrapped asset on the destination; the return leg
    \texttt{CrossBurnCall} burns wrapped and unlocks the escrow.
\end{itemize}
A single OHM transaction may carry at most one cross-chain action (send, recv,
transfer, or burn) and may not combine a cross action with an \ovm{} action.

\section{Express Bus and the Bridge}
\label{sec:express-bus}

\subsection{Express Bus}

The Express Bus is a per-route In/Out tracker for value flowing between \ox{}
and external systems.  Each bus $\beta$ tracks a \PriceAsset{} and accumulates
\[
f_\beta = \omega\cdot(\mathit{in} - \mathit{out})
\]
over a sampling window of $P$ blocks, where $\omega$ is a fixed-point weight.
A \emph{deposit} (master-approved, $\ge\texttt{quorum}$ of masters) tags
existing Spendable UTXOs as \BusReceipt{} outputs; a \emph{withdrawal}
(node-approved, $\ge\texttt{quorum}$ of servers) spends them and is additionally
owner-signed.  The signed flow signal is folded into the value-index formula
(Section~\ref{sec:value-formula}) so sustained outflow raises $V$ and
inflow lowers it.

\subsection{Bridge Module}

The bridge module provides verifiable cross-chain asset transfer among external
chains.  It maintains:
\begin{itemize}[leftmargin=2em]
  \item a \ConsensusProof{} (validator set, signer count, threshold, aggregated
    signature, epoch) for finality;
  \item a \ProofManager{} for Merkle-proof creation, validation, and caching;
  \item a \VerificationCoordinator{} managing verification sessions and
    reputation;
  \item a \SecureCommunicationLayer{} using X25519 key exchange with HKDF
    derivation and AES-256-GCM or ChaCha20-Poly1305 encryption at AES-128-CBC fallback;
  \item \emph{Recovery} logic for timed-out or failed sessions.
\end{itemize}
External \ChainId{}s are $64$-bit; proofs are verified against the source
chain's state commitment with a $2/3$ threshold signature.

\section{Network and Wire Protocol}
\label{sec:net}

\subsection{Transport}

Nodes communicate over a length-framed TCP transport.  Each frame is prefixed
by its size (little-endian) and capped at $16\,\mathrm{MiB}$.  Messages are
serialized deterministically and dispatched on the receiving end.

\subsection{Message Set}

The wire \Message{} enum covers liveness (\texttt{Ping}/\texttt{Pong}), chain
handshake, admission (node--master credential exchange), master governance
(proposal and epoch-transition broadcasts), consensus (proposal, attestation,
vote, quorum-certificate, and master-check messages), and wallet-service
queries (balance, UTXO set, inclusion proofs, block and transaction lookup,
value-index queries, transaction and operation submission, DAO and OHM state
queries).

\subsection{Master Selection}

Nodes select a master to connect to by minimum round-trip time, maintaining a
rolling RTT estimate and preferring the nearest reachable master for wallet
queries while broadcasting consensus messages to the full node set.

\section{Ephemeral Verification Keys and Secure Channels}
\label{sec:evk}

\subsection{EVK Key Chain}

Each block carries an optional \emph{ephemeral verification key}: an Ed25519
public key $\pi_h$ generated for block $h$ and the \emph{previous commitment}
$c_{h-1}$ from the prior block.  The \emph{chaining commitment} is
\[
c_h = \hash_{\mathrm{many}}\bigl(
  \texttt{``EVK\_CHAIN''}\,\|\,c_{h-1}\,\|\,\pi_h\,\|\,\mathrm{LE64}(h)\,\|\,\hash(B_h)
\bigr).
\]
The header stores $\pi_h$, $c_{h-1}$, and (optionally) $c_h$.  Masters and
auditors that do not retain private key material can validate the chain of
commitments from the genesis key ($\emptyroot$).

\subsection{Migration}

An optional \emph{migration} period makes the EVK extension mandatory-but-optional:
nodes may operate without keys while the fraction of keyed blocks is below
$\texttt{migration\_threshold\_percent}$ (default $90\%$).  Once the threshold
is met (or rolled back by an admin-signed transition), keys become mandatory for
all subsequently proposed blocks.  Checkpoints at fixed intervals
($\texttt{checkpoint\_interval}$) record commitment snapshots for efficient
catch-up verification.

\subsection{Secure Node Channels}

Authenticated, confidential peer channels are established via X25519 static key
exchange with an Ed25519 identity transcript binding.  Messages are AES-256-GCM
encrypt-then-MAC'd with a per-message sequence number and session identifier as
associated data; replays are rejected by a per-session sequence-set.  Private key
material is held in mlock'd, zeroizing buffers and erased on session completion.

\subsection{EVFS (File System Labeling)}
\label{sec:evfs}

The EVK chain extends to a content-addressable file system:
\begin{itemize}[label=--,leftmargin=2em]
  \item A \StorageLabel{} is a self-signed declaration of a node's capacity,
    used/unavailable bytes, and optional file metadata; v1 and v2 commitment
    domains are supported during migration.
  \item A \FileEvkLabel{} is a signed, chained per-file attestation binding
    content hash, size, and a ``EVFS\_FILE\_CHAIN'' commitment that links to the
    previous file, enabling stateless integrity verification without storing
    file bodies on-chain.
  \item Per-node file counts are capped; an epoch storage audit records
    compliance; deviations beyond $\texttt{storage\_label\_max\_deviation\_bps}$
    are rejected.
\end{itemize}
Max file size is $8\,\mathrm{MiB}$; max files per node is $10^5$.

\section{Rewards and Departures}
\label{sec:rewards}

\subsection{Block Rewards and Fees}

Each block's fees are split $4\%$ to masters and $96\%$ to server nodes by
default (\texttt{fee\_master\_bps}).  Minted block rewards are distributed
participation-weighted among servers (with a proposer bonus) and equally among
masters (derived from the server-reward derivative).  All earnings are
\LockedReward{} outputs, unlockable only through an approved withdrawal.

\subsection{Node Departure}
\label{sec:departure}

On a verified node departure (graceful or liveness-triggered), a
\texttt{DepartureEarnings} operation redistributes half of the departed
node's held earnings equally across connected members of both the server and
master groups; the confiscation fraction is
$\texttt{departure\_take\_bps}$ (default $50\%$, adjustable).  The remainder
stays locked for the departing node subject to the normal withdrawal path and
cooling-off period.

\section{Security and Invariants}
\label{sec:security}
\label{sec:analysis}

The protocol's safety rests on the following invariants, each enforced by the
master audit and/or server replay:

\begin{invariant}[Conservation]
For every transaction, $\sum \mathrm{inputs} = \sum \mathrm{outputs} + \mathrm{fee}$,
and all inputs exist, are owned by the spender, and (if rebound) are mature.
\end{invariant}

\begin{invariant}[Lineage]
$\sum_{\mathit{live\;UTXO}} \mathit{amount} = \Sigma_h$ at every block $h$.
\end{invariant}

\begin{invariant}[Header consistency]
$R_T, R_O, R_{\mathcal{R}}$, the value index, generation, and supply in the header
must be independently recomputable from the block body and prior state.
\end{invariant}

\begin{invariant}[Governance threshold]
No mint, burn, withdraw, price change, epoch transition, or node admission
takes effect without the configured super-majority of the appropriate
governing set, verified against the current epoch's GovernanceSet.
\end{invariant}

\begin{invariant}[Determinism]
All on-chain arithmetic is integer fixed-point; no floating point, no
platform-dependent operations, and no nondeterministic control flow exist in the
state transition or \ovm{} execution.
\end{invariant}

\subsection{Treatment of the Value Index}
\label{sec:bunnyhop}

The value index $V$ is a monotonically tracked scalar in $[0, \infty)$ that
ratchets downward on exit pressure and upward on activity, smoothed by
$\lambda$.  When pressure is severe enough to drive the raw value non-positive,
the protocol \emph{rebounds}: it pins $V$ to the last positive value, increments
the generation counter, and marks exit-reserve credits as refundable.  Rebound
credits are subject to bunny-hop maturity (Section~\ref{sec:state}) to prevent
rapid re-spending of the refund.

\section{Related Subsystems}

\subsection{Timer}
A betting subsystem supporting timestamp-direction wagers (up/down) with an
accumulation period, configurable min/max bet amounts, and resolution at
expiration.  Correct-direction bettors share the losing side's stake.

\subsection{Genesis and Parameters}
The genesis document allocates the full $S_0$ to recipient addresses and
declares the initial master and server-node sets, the \texttt{CodeManifest}
approval registry, and the \ChainParams{} block.  All parameters are rational
or fixed-point integers chosen so that every node derives identical results.

\section{Conclusion}

\ox{} presents a layered architecture in which a small auditor set can
transitively trust the full-state chain through Merkle proofs, deterministic
replay, and supply-conservation checks, while a rich execution environment
(\ovm{}), tenant chains (\ohms{}), fuel currencies (\gas{}, Speedometer),
governance (DAO), and cross-chain plumbing (Express Bus, Bridge) are built
\emph{as state machine extensions} rather than as separate ledgers.  The
integration of an integer fixed-point value index with bunny-hop rebounds,
threshold-governed supply changes, and ephemeral per-block verification keys
yields a self-contained monetary and consensus design that is fully specified
by the block history and the genesis parameters.

\bibliographystyle{plain}
\bibliography{oxydome}

\end{document}
